\documentclass[11pt,a4paper]{article}

\usepackage{fontspec}
\IfFontExistsTF{Times New Roman}{\setmainfont{Times New Roman}}{\setmainfont{Tinos}}
\IfFontExistsTF{Consolas}{\setmonofont{Consolas}[Scale=0.92]}{\setmonofont{Liberation Mono}[Scale=0.92]}
\usepackage[english]{babel}
\usepackage{geometry}
\usepackage{xcolor}
\usepackage{enumitem}
\usepackage{listings}
\usepackage{booktabs}
\usepackage{array}
\usepackage{amsmath,amssymb}
\usepackage{titlesec}
\usepackage{fancyhdr}
\usepackage{parskip}
\usepackage{tcolorbox}
\tcbuselibrary{skins}
\usepackage{tabularx}
\usepackage{lastpage}
\usepackage{hyperref}

\geometry{top=1.5cm,bottom=1.5cm,left=1.7cm,right=1.7cm}

\definecolor{primaryblue}{RGB}{18,65,115}
\definecolor{secondaryteal}{RGB}{14,102,110}
\definecolor{SoftBlue}{RGB}{232,242,250}
\definecolor{CodeBg}{RGB}{248,250,252}
\definecolor{DarkGray}{RGB}{55,55,55}
\definecolor{boxbg}{RGB}{248,250,252}
\definecolor{borderblue}{RGB}{150,185,215}

\setlength{\parindent}{0pt}
\setlength{\parskip}{0pt}
\setlist[itemize]{leftmargin=1.3em,itemsep=1pt,topsep=1pt,parsep=0pt,partopsep=0pt}
\setlist[enumerate]{leftmargin=1.5em,itemsep=1.5pt,topsep=1pt,parsep=0pt,partopsep=0pt}
\setlist[enumerate,2]{topsep=0.5pt,itemsep=0pt,parsep=0pt,partopsep=0pt}

\lstset{
  language=Python,
  basicstyle=\ttfamily\scriptsize,
  keywordstyle=\color{blue!75!black}\bfseries,
  stringstyle=\color{red!65!black},
  commentstyle=\color{primaryblue!80!black}\itshape,
  showstringspaces=false,
  columns=fullflexible,
  keepspaces=true,
  breaklines=true,
  frame=single,
  rulecolor=\color{primaryblue!45},
  backgroundcolor=\color{CodeBg},
  xleftmargin=0.2em,
  xrightmargin=0.2em,
  aboveskip=0.2em,
  belowskip=0.2em,
  tabsize=4
}

\pagestyle{fancy}
\fancyhf{}
\setlength{\headheight}{14pt}
\fancyhead[L]{\footnotesize\textbf{DSAI1003 -- Fundamental Programming Concepts in Python}}
\fancyhead[R]{\footnotesize\textbf{Loops \& Iteration Patterns Exam -- 60 Minutes}}
\fancyfoot[L]{\scriptsize Faculty of Data Science and Artificial Intelligence -- National Economics University}
\fancyfoot[R]{\footnotesize Page \thepage\ of 6}
\renewcommand{\headrulewidth}{0.4pt}
\renewcommand{\footrulewidth}{0.3pt}

\titleformat{\section}{\large\bfseries\color{primaryblue}}{}{0pt}{}
\titleformat{\subsection}{\normalsize\bfseries\color{primaryblue}}{}{0pt}{}
\titlespacing*{\section}{0pt}{3.5pt}{1.5pt}
\titlespacing*{\subsection}{0pt}{2.5pt}{1pt}

\newtcolorbox{headerbox}{
  colback=white,
  colframe=primaryblue,
  boxrule=0.7pt,
  arc=1.5mm,
  left=2.5mm,
  right=2.5mm,
  top=1.2mm,
  bottom=1.2mm
}

\newtcolorbox{answerbox}[2][]{
  enhanced,
  colback=white,
  colframe=DarkGray!45,
  boxrule=0.5pt,
  arc=1mm,
  left=2mm,
  right=2mm,
  top=1.2mm,
  bottom=1.2mm,
  title=\footnotesize\color{DarkGray}\textbf{#2},
  #1
}

\hypersetup{
  colorlinks=true,
  linkcolor=primaryblue,
  urlcolor=primaryblue,
  pdftitle={DSAI1003 - Loops and Iteration Patterns in Python Exam},
  pdfauthor={Faculty of Data Science and Artificial Intelligence - National Economics University (NEU)}
}

\begin{document}

% =========================================================================
% PAGE 1: OFFICIAL HEADER + ANSWER SHEET + SECTION A (Q1 - Q5)
% =========================================================================

\noindent
\begin{minipage}[t]{0.64\textwidth}
  \raggedright\small
  \textbf{NATIONAL ECONOMICS UNIVERSITY}\\
  \textbf{COLLEGE OF TECHNOLOGY}\\
  \textbf{FACULTY OF DATA SCIENCE \& ARTIFICIAL INTELLIGENCE (FDA)}
\end{minipage}%
\hfill
\begin{minipage}[t]{0.34\textwidth}
  \raggedleft\footnotesize
  \textbf{OFFICIAL PROGRESS EXAMINATION}\\
  Academic Year: 2026--2027\\
  Semester: Fall 2026 \quad Course Code: \textbf{DSAI1003}
\end{minipage}

\vspace{0.15em}

\begin{center}
  {\Large\bfseries\color{primaryblue} FUNDAMENTAL PROGRAMMING CONCEPTS IN PYTHON}\\[0.08em]
  {\normalsize\bfseries PROGRESS EXAM: LOOPS \& ITERATION PATTERNS}\\[0.08em]
  {\normalsize\textbf{Paper Code:} \rule{1.6cm}{0.4pt} \quad | \quad \textbf{Duration: 60 Minutes} (Excluding paper distribution)}
\end{center}

\vspace{0.05em}

\begin{headerbox}
\begin{tabular*}{\textwidth}{@{\extracolsep{\fill}}ll@{}}
  \textbf{Candidate Full Name:} \dotfill & \textbf{Student ID (MSSV):} \dotfill \\[0.2em]
  \textbf{Class / Cohort:} \dotfill & \textbf{Exam Room / Seat No.:} \dotfill
\end{tabular*}

\vspace{0.2em}
\renewcommand{\arraystretch}{1.05}
\setlength{\tabcolsep}{4.5pt}
\centering
\begin{tabular}{|l|c|c|c|c||c|c|}
\hline
\textbf{Section} & \textbf{A} & \textbf{B} & \textbf{C} & \textbf{D} & \textbf{Total (100)} & \textbf{Examiner Signature} \\
\hline
\textbf{Max Score} & 20 & 25 & 25 & 30 & \textbf{100} & \\[0.3em]
\cline{1-6}
\textbf{Score} & & & & & & \\[0.4em]
\hline
\end{tabular}

\vspace{0.15em}
\raggedright
\scriptsize
\textbf{Instructions:} Closed-book examination. All electronic devices and internet access are prohibited. Answer all questions directly in the designated spaces. In programming tasks, apply iteration structures (\texttt{while}, \texttt{for}), \texttt{range()}, sentinel handling, accumulator patterns, and nested loops. Do \textbf{not} define custom functions (\texttt{def}) unless explicitly requested. Assume user input follows the requested type format. Do not detach pages.
\end{headerbox}

\vspace{1mm}
\noindent
\textbf{MULTIPLE-CHOICE ANSWER SHEET} \textit{(Write A, B, C, or D in the corresponding box)}:
\begin{center}
\setlength{\tabcolsep}{5pt}
\renewcommand{\arraystretch}{1.0}
\begin{tabular}{|c|c|c|c|c|c|c|c|c|c|c|}
\hline
\textbf{Question} & \textbf{1} & \textbf{2} & \textbf{3} & \textbf{4} & \textbf{5} & \textbf{6} & \textbf{7} & \textbf{8} & \textbf{9} & \textbf{10} \\
\hline
\textbf{Answer} & & & & & & & & & & \\[1.8mm]
\hline
\end{tabular}
\end{center}
\vspace{-2.5mm}
\section*{Section A: Loops \& Iteration Concepts \hfill \normalsize\normalfont(20 points)}
\textit{Answer all 10 questions. Each correct answer carries 2 points. Suggested time: 12 minutes.}

\begin{enumerate}[label=\textbf{Q\arabic*.},leftmargin=2em,itemsep=1.5pt,topsep=1pt]

\item When should a programmer select an event-controlled \texttt{while} loop rather than a count-controlled \texttt{for} loop?
\begin{enumerate}[label=\Alph*.,leftmargin=2em]
  \item When the exact number of iterations is known in advance before entering the loop.
  \item When repeating a block until an external event or sentinel condition occurs during execution.
  \item When traversing every character of a fixed-length string from start to finish.
  \item When generating an arithmetic progression of integers with a constant step size.
\end{enumerate}

\item What sequence of integers is generated by the expression \texttt{list(range(10, 2, -3))}?
\begin{enumerate}[label=\Alph*.,leftmargin=2em]
  \item \texttt{[10, 7, 4]}
  \item \texttt{[10, 7, 4, 1]}
  \item \texttt{[10, 8, 6, 4, 2]}
  \item An empty list \texttt{[]}, because the starting value is greater than the stopping value.
\end{enumerate}

\item What is the primary cause of an \textbf{infinite loop} in a \texttt{while} statement?
\begin{enumerate}[label=\Alph*.,leftmargin=2em]
  \item The loop body executes too quickly for the CPU register cache to update.
  \item The continuation condition never evaluates to \texttt{False} because the control variable is not modified.
  \item The loop counter exceeds Python's maximum integer limit (\texttt{sys.maxsize}).
  \item The loop contains an \texttt{if} statement with more than three \texttt{elif} branches.
\end{enumerate}

\item An algorithm intended to print numbers from 1 to 5 inclusive prints 1 to 4 because the condition is \texttt{count < 5} instead of \texttt{count <= 5}. What type of defect is this?
\begin{enumerate}[label=\Alph*.,leftmargin=2em]
  \item A syntax error caught during byte-code compilation.
  \item A name collision error in local namespace resolution.
  \item An off-by-one error, which is a classic form of logic error.
  \item A runtime exception raising an \texttt{IndexError}.
\end{enumerate}

\item Why is a \textbf{priming read} essential in a sentinel-controlled \texttt{while} loop?
\begin{enumerate}[label=\Alph*.,leftmargin=2em]
  \item It converts user input into a floating-point value automatically.
  \item It obtains the initial input value before testing the condition, preventing the sentinel from being processed as valid data.
  \item It forces Python's garbage collector to allocate memory for the accumulator variable.
  \item It causes the loop to run at least twice regardless of the input value.
\end{enumerate}

\end{enumerate}

\newpage

% ==============================================================================
% PAGE 2: SECTION A (Q6 - Q10) + SECTION B Q11
% ==============================================================================

\begin{enumerate}[label=\textbf{Q\arabic*.},leftmargin=2em,itemsep=1.5pt,topsep=1pt,resume]

\item In a nested loop where the outer loop iterates $R$ times and the inner loop iterates $C$ times, how many total times does the inner loop body execute?
\begin{enumerate}[label=\Alph*.,leftmargin=2em]
  \item $R + C$ times
  \item $R \times C$ times
  \item $R^C$ times
  \item Exactly $\max(R, C)$ times
\end{enumerate}

\item What is the operational distinction between the \texttt{break} and \texttt{continue} statements?
\begin{enumerate}[label=\Alph*.,leftmargin=2em]
  \item \texttt{break} restarts the current iteration from the beginning; \texttt{continue} terminates the entire program.
  \item \texttt{break} immediately exits the enclosing loop; \texttt{continue} skips the remainder of the current iteration and advances to the next.
  \item \texttt{break} is valid only in \texttt{for} loops; \texttt{continue} is valid only in \texttt{while} loops.
  \item \texttt{break} skips the inner loop only; \texttt{continue} exits both inner and outer loops simultaneously.
\end{enumerate}

\item When searching for the maximum value in a non-empty sequence of positive and negative numbers, what is the safest initialization strategy for \texttt{max\_val}?
\begin{enumerate}[label=\Alph*.,leftmargin=2em]
  \item Always initialize \texttt{max\_val = 0}.
  \item Initialize \texttt{max\_val} to the first element of the sequence or negative infinity (\texttt{float('-inf')}).
  \item Always initialize \texttt{max\_val = 999999}.
  \item Leave \texttt{max\_val} unassigned until the first comparison statement executes.
\end{enumerate}

\item What does the following loop output: \texttt{for ch in "NEU2026": if ch.isdigit(): print(ch, end="")}?
\begin{enumerate}[label=\Alph*.,leftmargin=2em]
  \item \texttt{NEU}
  \item \texttt{2026}
  \item \texttt{N E U 2 0 2 6}
  \item \texttt{True}
\end{enumerate}

\item In Python, when does the optional \texttt{else} block attached to a \texttt{for} or \texttt{while} loop execute?
\begin{enumerate}[label=\Alph*.,leftmargin=2em]
  \item Every time an iteration finishes normally.
  \item Only when the loop encounters a \texttt{break} statement.
  \item When the loop completes all scheduled iterations without encountering a \texttt{break}.
  \item Only when an unexpected runtime exception is raised inside the loop.
\end{enumerate}

\end{enumerate}

\vspace{1mm}
\section*{Section B: Output Tracing, Logic Diagnosis \& Error Correction \hfill \normalsize\normalfont(25 points)}
\textit{Answer all questions Q11--Q15 directly in the designated spaces. Suggested time: 15 minutes.}

\subsection*{Q11. Exact Output Tracing with Nested Loops, Accumulator \& Conditions (5 points)}
Trace the following program carefully and write its exact console output:

\begin{lstlisting}
total = 0
for i in range(1, 4):
    for j in range(1, 4):
        if i == j:
            continue
        if i + j > 4:
            break
        total += i * 10 + j
        print(f"[{i},{j}]", end=" ")
print(f"\nFinal Total: {total}")
\end{lstlisting}

\begin{answerbox}[height=3.4cm]{Output Box for Q11}
\end{answerbox}

\newpage

% ==============================================================================
% PAGE 3: SECTION B, Q12--Q15
% ==============================================================================

\subsection*{Q12. Complete Variable State Hand-Trace Table (5 points)}
Trace the following \texttt{while} loop until termination. Fill in the state trace table for each iteration:

\begin{minipage}[t]{0.48\textwidth}
\begin{lstlisting}
n = 432
rev = 0
step = 0
while n > 0:
    digit = n % 10
    rev = rev * 10 + digit
    n = n // 10
    step += 1
\end{lstlisting}
\end{minipage}%
\hfill
\begin{minipage}[t]{0.48\textwidth}
\centering
\small
\begin{tabular}{|c|c|c|c|c|}
\hline
\textbf{Step} & \textbf{\texttt{n}} & \textbf{\texttt{digit}} & \textbf{\texttt{rev}} & \textbf{\texttt{n > 0}} \\
\hline
Initial & 432 & -- & 0 & True \\[1.5mm]
\hline
1 & & & & \\[1.5mm]
\hline
2 & & & & \\[1.5mm]
\hline
3 & & & & \\[1.5mm]
\hline
Final & & -- & & False \\[1.5mm]
\hline
\end{tabular}
\end{minipage}

\vspace{1mm}
\subsection*{Q13. Diagnose and Correct Three Common Loop Defects (5 points)}
Identify the exact error and write the corrected code snippet in the table below:

\renewcommand{\arraystretch}{1.15}
\begin{tabularx}{\textwidth}{|X|p{4.2cm}|p{6.0cm}|}
\hline
\textbf{Flawed Snippet} & \textbf{Defect Classification} & \textbf{Exact Corrected Line(s)} \\
\hline
\texttt{i = 1; sum\_val = 0}\\
\texttt{while i <= 10:}\\
\texttt{\ \ \ \ sum\_val += i}\\
(Intended: sum 1 to 10) & & \\[3.2mm]
\hline
\texttt{total = 0}\\
\texttt{for i in range(1, 5):}\\
\texttt{\ \ \ \ total += i}\\
(Intended: sum 1 to 5 inclusive) & & \\[3.2mm]
\hline
\texttt{for i in range(3):}\\
\texttt{\ \ \ \ val = input("Enter: ")}\\
\texttt{print("Read:", val)}\\
(Intended: print every entry) & & \\[3.2mm]
\hline
\end{tabularx}

\vspace{1mm}
\subsection*{Q14. String Iteration \& Filtering Accumulator Pattern (5 points)}
Predict the exact printed output of the following character aggregation program:

\begin{lstlisting}
text = "DSFE-Ai@2026#VN"
digits = ""
vowels = ""
vowel_set = "aeiouAEIOU"
for ch in text:
    if ch.isdigit():
        digits += ch
    elif ch in vowel_set:
        vowels += ch.upper()
print(f"Vowels: {vowels}, Digits: {digits}")
print(f"Total Kept: {len(vowels) + len(digits)}")
\end{lstlisting}

\begin{answerbox}[height=2.2cm]{Output Box for Q14}
\end{answerbox}

\vspace{1mm}
\subsection*{Q15. Sentinel Loop Architecture \& Protection Against Anomalies (5 points)}
\begin{itemize}[leftmargin=2em]
  \item Explain why the statement \texttt{total += val} placed immediately after \texttt{val = float(input())} in a \texttt{while val != -1:} loop introduces a logic defect when the sentinel \texttt{-1} is entered: \hrulefill
  \item What conditional check must be placed \textbf{before} dividing by \texttt{count} when computing the average of sentinel-entered numbers? \hrulefill
\end{itemize}

\newpage

% ==============================================================================
% PAGE 4: SECTION C
% ==============================================================================

\section*{Section C: Program Design \& Logic Specification \hfill \normalsize\normalfont(25 points)}
\textit{Answer all tasks. Suggested time: 15 minutes.}

\begin{tcolorbox}[colback=boxbg,colframe=secondaryteal,arc=1.5mm,boxrule=0.6pt,top=1.5mm,bottom=1.5mm]
\textbf{Problem Specification: Retail Cash Register \& Daily Sales Aggregator}\\[0.5mm]
\small
A retail shop requires an automated sales transaction aggregator that processes invoice amounts until the cashier signals end-of-shift by entering \texttt{-1} (sentinel):
\begin{itemize}[itemsep=0.5pt,topsep=0.5pt]
  \item \textbf{Input Validation:} Any negative transaction other than \texttt{-1} is invalid. The program must warn the cashier and prompt again without adding to the transaction count or revenue. Zero is also rejected.
  \item \textbf{Accumulation:} Keep track of: (1) \texttt{total\_revenue}, (2) \texttt{valid\_count}, (3) \texttt{max\_sale}, and (4) \texttt{high\_value\_count} (sales $\ge \$500.00$).
  \item \textbf{Summary Metrics:} At termination, calculate the \texttt{average\_sale}. If no valid transactions were recorded, display an appropriate notice and set all monetary metrics to \$0.00.
  \item \textbf{VIP Flag:} If \texttt{total\_revenue} exceeds \$2,000.00 or \texttt{high\_value\_count} $\ge 3$, flag the shift as \texttt{"Target Exceeded"}.
\end{itemize}
\end{tcolorbox}

\vspace{0.5mm}
\textbf{Task 1: Identify Loop Type, Sentinel \& State Variables} (6 points)
\begin{itemize}[leftmargin=2em]
  \item \textbf{Loop Structure \& Termination Condition:} \hrulefill\\[3mm]
  \item \textbf{Accumulator \& Counter Variables:} \hrulefill\\[3mm]
  \item \textbf{Initialization Values for State Variables:} \hrulefill
\end{itemize}

\vspace{0.5mm}
\textbf{Task 2: Write the Structured Iteration Logic in Ordered Steps} (7 points)\\
Formulate sequential pseudo-logic showing the priming read, loop condition, validation branch, accumulators, and post-loop summary.

\begin{answerbox}[height=5.2cm]{Program-Logic Work Area}
\end{answerbox}

\vspace{0.5mm}
\textbf{Task 3: Hand Trace / Step-by-Step Test Case Execution} (6 points)\\
Given the input stream: \texttt{150.0}, \texttt{-25.0} (invalid), \texttt{600.0}, \texttt{350.0}, \texttt{-1.0} (sentinel). Trace the final metrics:

\begin{answerbox}[height=3.5cm]{Calculation Box for Task 3}
\end{answerbox}

\vspace{0.5mm}
\textbf{Task 4: Explain Two Design Decisions} (6 points)
\begin{enumerate}[label=\alph*),leftmargin=2em]
  \item Why is a \texttt{while} loop strictly required instead of a \texttt{for} loop for this cashier application?\\[3.5mm]
  \hrulefill
  \item How does your logic prevent a \texttt{ZeroDivisionError} if the cashier enters \texttt{-1} on the very first prompt?\\[3.5mm]
  \hrulefill
\end{enumerate}

\newpage

% ==============================================================================
% PAGE 5: SECTION D, PROBLEM 1
% ==============================================================================

\section*{Section D: Applied Programming in Python \hfill \normalsize\normalfont(30 points)}
\textit{Write complete, syntactically correct Python statements. Apply proper iteration structures (\texttt{while}, \texttt{for}), input validation, and formatting. Do \textbf{not} define custom functions (\texttt{def}). Suggested time: 18 minutes.}

\subsection*{Problem 1: Compound Savings Target \& Milestone Growth Simulator (15 points)}
Write a Python program that simulates the year-by-year compounding growth of an investment until it reaches or exceeds a specified target amount.

\textbf{Requirements:}
\begin{enumerate}[leftmargin=2em,label=\arabic*.,itemsep=1pt]
  \item Prompt for initial deposit \texttt{principal} (float) and annual interest rate \texttt{annual\_rate} in percent (float).
  \item Prompt for the financial goal \texttt{target\_balance} (float).
  \item \textbf{Input Validation:} Use a \texttt{while} loop to ensure \texttt{principal > 0}, \texttt{annual\_rate > 0}, and \texttt{target\_balance > principal}. Re-prompt until valid.
  \item Use a \texttt{while} loop to simulate growth year by year until \texttt{balance >= target\_balance}.
  \item In each year: \texttt{interest = balance * (annual\_rate / 100)}; \texttt{balance += interest}.
  \item Print a formatted progress table showing: \texttt{Year}, \texttt{Start Balance}, \texttt{Interest Earned}, and \texttt{End Balance} (all balances formatted with 2 decimal places).
  \item After loop termination, display: Total Years Required, Final Balance, and Total Interest Earned across all years.
\end{enumerate}

The output should follow this format:
\begin{lstlisting}
Year    Start Balance     Interest Earned   End Balance      
-------------------------------------------------------------
1       $1000.00          $70.00            $1070.00         
2       $1070.00          $74.90            $1144.90         
...
Target of $2000.00 reached in 11 years! Total Interest: $1104.89
\end{lstlisting}

\begin{answerbox}[height=13.2cm]{Python Solution for Problem 1}
\end{answerbox}

\newpage

% ==============================================================================
% PAGE 6: SECTION D, PROBLEM 2
% ==============================================================================

\subsection*{Problem 2: Numerical Sequence Scanner \& Adjacent Fluctuation Analyzer (15 points)}
Write a Python program that scans a sequence of daily stock price quotations, computes statistical summary values, and monitors adjacent day-to-day market fluctuations.

\textbf{Requirements:}
\begin{enumerate}[leftmargin=2em,label=\arabic*.,itemsep=1pt]
  \item Prompt the user for the number of trading days $N$ (integer). Validate with a loop so that $N \ge 3$.
  \item Using a \texttt{for} loop running $N$ times, prompt the user for each day's price (float, $> 0$).
  \item Accumulate the sum to compute the overall \texttt{mean\_price} after all inputs are read.
  \item Track the \textbf{highest price} and \textbf{lowest price} observed across the $N$ days.
  \item For every day after Day 1, compute the price change from the previous day: $\Delta = \text{price}_{\text{current}} - \text{price}_{\text{previous}}$.
  \item Track: (a) count of gain days ($\Delta > 0$), (b) count of loss days ($\Delta < 0$), and (c) the \textbf{largest single-day surge} ($\max \Delta$).
  \item Display a neatly formatted report showing: Number of Days, Average Price, Highest Price, Lowest Price, Days Up, Days Down, and Maximum Single-Day Surge (all prices formatted to 2 decimal places).
\end{enumerate}

The output should follow this format:
\begin{lstlisting}
--- MARKET SEQUENCE ANALYSIS REPORT ---
Trading Days Analyzed: 5
Price Range: $102.50 to $118.20 | Mean Price: $110.84
Upward Days: 3 | Downward Days: 1
Largest Single-Day Surge: +$8.50
\end{lstlisting}

\begin{answerbox}[height=15.2cm]{Python Solution for Problem 2}
\end{answerbox}

\vfill
\begin{center}
  \rule{0.5\textwidth}{0.4pt}\\[1mm]
  \textbf{\color{primaryblue}--- END OF EXAM PAPER ---}
\end{center}

\end{document}
